% Part: counterfactuals
% Chapter: introduction
% Section: strict-conditional

\documentclass[../../../include/open-logic-section]{subfiles}

\begin{document}

\olfileid{cnt}{int}{str}

\olsection{سخت شرطیہ}

لوئس نے \emph{سخت شرطیہ}~$\strictif$ پیش کیتا تے دلیل دتی کہ استلزام
ایہدے نال میل کھاندا اے، مادی شرطیے نال نہیں۔ سچائی دی امکانی منطق وچ
$!A \strictif !B$ نوں $\Box(!A \lif !B)$ دے طور اُتے تعریف کر سکدے آں۔
سو سخت شرطیہ (کسے جہان اُتے) اُدوں تے صرف اُدوں سچ اے جدوں اوہدے نال
ملدا مادی شرطیہ لازمی اے۔

مادی شرطیے دے معمّیاں دے لحاظ نال سخت شرطیہ کیویں چلدا اے؟ جھوٹے مقدم
والا سخت شرطیہ تے سچی تالی والا سخت شرطیہ، دونیں سچ وی ہو سکدے نیں
تے جھوٹ وی۔ نالے $(!A \strictif !B) \lor (!B \strictif !A)$ معتبر نہیں۔
سخت شرطیہ $!A \strictif !B$، $\lnot !A \lor !B$ دے برابر وی نہیں،
سو ایہ صداقتی-فنکشنی نہیں۔

ساڈے کول ایہ نیں:
\begin{align}
  !A \strictif !B & \Entails \lnot !A \lor !B
  \text{ پر }\\
  \lnot !A \lor !B & \Entails/ !A \strictif !B\\
  !B & \Entails/ !A \strictif !B\\
  \lnot !A & \Entails/ !A \strictif !B\\
  \lnot(!A \lif !B) & \Entails/ !A \land \lnot !B
  \text{ پر }\\
  !A \land \lnot !B & \Entails \lnot(!A \strictif !B)
  \intertext{فیر وی سخت شرطیے لئی اثباتِ مقدم دا قاعدہ چلدا اے:}
  !A, !A \strictif !B & \Entails !B
  \intertext{سخت شرطیاں دیاں تالیاں اکٹھیاں ہو سکدیاں نیں:}
  !A \strictif !B,  !A \strictif !C & \Entails !A \strictif (!B \land !C)
  \intertext{سخت شرطیے لئی مقدم نوں مضبوط کرن دا قاعدہ چلدا اے:}
  !A \strictif !B & \Entails (!A \land !C) \strictif !B
  \intertext{سخت شرطیہ متعدی وی اے:}
  !A \strictif !B, !B \strictif !C & \Entails !A \strictif !C
  \intertext{آخر وچ، سخت شرطیہ اپنے عکسِ نقیض دے ہم ارز اے:}
  !A \strictif !B & \Entails \lnot !B \strictif \lnot !A\\
  \lnot !B \strictif \lnot !A & \Entails !A \strictif !B
\end{align}

\begin{prob}
  استلزام دے اوہناں تعلقاں لئی \Log{S5} وچ ردّی مثالاں دیو جیہڑے سخت
  شرطیے لئی قائم نہیں ہوندے، یعنی ہِیٹھلیاں لئی:
  \begin{enumerate}
  \item $\lnot p \Entails/ \Box(p \lif q)$
  \item $q \Entails/ \Box(p \lif q)$
  \item $\lnot\Box(p \lif q) \Entails/ p \land \lnot q$
  \item $\Entails/ \Box(p \lif q) \lor \Box(q \lif p)$
  \end{enumerate}
\end{prob}

\begin{prob}
  ہِیٹھلیاں گلاں دے \Log{S5} وچ ثبوت دے کے وکھاؤ کہ معتبر استلزام دے
  تعلق سخت شرطیے لئی قائم نیں:
  \begin{enumerate}
  \item  $\Box(!A \lif !B) \Entails \lnot !A \lor !B$
  \item $!A \land \lnot !B \Entails \lnot\Box(!A \lif !B)$
  \item $!A, \Box(!A \lif !B) \Entails !B$
  \item $\Box(!A \lif !B), \Box(!A \lif !C) \Entails \Box(!A \lif (!B
    \land !C))$
  \item $\Box(!A \lif !B) \Entails \Box((!A \land !C) \lif !B)$
  \item $\Box(!A \lif !B), \Box(!B \lif !C) \Entails \Box(!A
    \lif !C)$
  \item $\Box(!A \lif !B) \Entails \Box(\lnot !B \lif \lnot !A)$
  \item $\Box(\lnot !B \lif \lnot !A) \Entails \Box(!A \lif !B)$
  \end{enumerate}
\end{prob}

پر سخت شرطیے دے اپنے وی ’معمّے‘ نیں۔ جیویں جھوٹے مقدم یا سچی تالی
والا مادی شرطیہ سچ ہوندا اے، اویویں \emph{لازمی طور اُتے} جھوٹے مقدم یا
لازمی طور اُتے سچی تالی والا سخت شرطیہ سچ ہوندا اے۔ نالے ہر سچا سخت شرطیہ
لازمی طور اُتے سچ اے، تے ہر جھوٹا سخت شرطیہ لازمی طور اُتے جھوٹ۔
دوجے لفظاں وچ،
\begin{align}
  \Box \lnot !A & \Entails !A \strictif !B\\
  \Box !B & \Entails !A \strictif !B\\
  !A \strictif !B& \Entails \Box(!A \strictif !B)\\
  \lnot(!A \strictif !B) & \Entails \Box\lnot(!A \strictif !B)
\end{align}
جے $\strictif$ دا مطلب ’لازم کردا اے‘ لئیے، تاں ایہ اوکڑاں نہیں۔
آخر منطقی استلزام دے تعلق ریاضیاتی حقیقتاں نیں، ایس لئی اوہ اتفاقی نہیں ہو سکدے۔
پر جے اسیں $\strictif$ نوں ’جے \dots تاں \dots‘ دا مطلب دین والا منطقی
رابطہ بنائیے، تاں ایہناں نال اوکڑ پَندی اے، خاص کر کے آخری دوناں نال۔
کیوں جے ’جے \dots تاں \dots‘ والے کجھ بیان اتفاقی طور اُتے سچ یا اتفاقی
طور اُتے جھوٹ ہوندے نیں۔ اصل وچ عام طور اُتے اوہ نہ لازمی ہوندے نیں،
نہ ناممکن۔

\begin{prob}
  ہِیٹھلیاں گلاں دے \Log{S5} وچ ثبوت دیو:
  \begin{enumerate}
    \item  $\Box \lnot !A \Entails !A \strictif !B$
    \item $!A \strictif !B \Entails \Box(!A \strictif !B)$
    \item $\lnot(!A \strictif !B)  \Entails \Box\lnot(!A \strictif !B)$
  \end{enumerate}
  ایس لئی $\strictif$ دی تعریف ورتو۔
\end{prob}

\end{document}
